\documentclass[11pt]{article}
\usepackage[a4paper,margin=1in]{geometry}
\usepackage{amsmath,amssymb,mathtools}
\usepackage{enumitem}
\usepackage{unicode-math}
\setlist{leftmargin=*}
\allowdisplaybreaks
\emergencystretch=3em
\title{KKT Normalization and Sonin Reduction}
\author{}
\date{}
\begin{document}
\maketitle
I use the KKT normalization
\begin{equation*}
g(x)=g_n^{(\alpha,\beta)}(x)
=C_{n,\alpha,\beta}
\left(\frac{1-x}{2}\right)^{\alpha/2}
\left(\frac{1+x}{2}\right)^{\beta/2}
P_n^{(\alpha,\beta)}(x),
\end{equation*}
where
\begin{equation*}
C_{n,\alpha,\beta}^2
=
\frac{\Gamma(n+1)\Gamma(n+\alpha+\beta+1)}
{\Gamma(n+\alpha+1)\Gamma(n+\beta+1)}.
\end{equation*}
This is the normalization used in KKT; they record the same \(g\), its \(L^2\)-norm, and the refined integer-parameter estimate corresponding to your target inequality.

Let
\begin{equation*}
A=1-x^2,\qquad s=\alpha+\beta,\qquad d=\beta-\alpha,
\end{equation*}
\begin{equation*}
U(x)=d-sx=\beta-\alpha-(\alpha+\beta)x,
\end{equation*}
\begin{equation*}
\lambda=n(n+s+1),\qquad \kappa=\lambda+\frac{s}{2}.
\end{equation*}

\section*{1. Exact transformed ODE}

The Jacobi equation is
\begin{equation*}
A P''+\bigl(d-(s+2)x\bigr)P'+\lambda P=0.
\end{equation*}
Writing
\begin{equation*}
g=e^\phi P,\qquad
\phi'(x)=\frac{U(x)}{2A(x)},
\end{equation*}
one obtains, after direct substitution,
\begin{equation}
\boxed{
g''-\frac{2x}{1-x^2}g'
+
\frac{\kappa(1-x^2)-\frac14\bigl(d-sx\bigr)^2}{(1-x^2)^2},g=0.
}
\tag{1}
\end{equation}
Equivalently, in Sturm–Liouville form,
\begin{equation}
\boxed{
\bigl(A g'\bigr)'
+
\left(
\kappa-\frac{U(x)^2}{4A}
\right)g=0.
}
\tag{2}
\end{equation}

Removing the first derivative in the same \(x\)-variable: set
\begin{equation*}
Y(x)=\sqrt{A(x)},g(x).
\end{equation*}
Then
\begin{equation}
\boxed{
Y''+Q_x(x)Y=0,
}
\tag{3}
\end{equation}
where
\begin{equation}
\boxed{
Q_x(x)
=
\frac{
1+\kappa(1-x^2)-\frac14(d-sx)^2
}
{(1-x^2)^2}.
}
\tag{4}
\end{equation}

The cleaner Liouville form uses \(x=\cos\theta\), \(0<\theta<\pi\), and
\begin{equation*}
Z(\theta)=\sqrt{\sin\theta},g(\cos\theta).
\end{equation*}
Then
\begin{equation}
\boxed{
Z''(\theta)+Q_\theta(\theta)Z(\theta)=0,
}
\tag{5}
\end{equation}
with
\begin{equation}
\boxed{
Q_\theta(\theta)
=
\left(n+\frac{\alpha+\beta+1}{2}\right)^2
+
\frac{\frac14-\alpha^2}{4\sin^2(\theta/2)}
+
\frac{\frac14-\beta^2}{4\cos^2(\theta/2)}.
}
\tag{6}
\end{equation}
This is the standard Jacobi Liouville normal form.

\section*{2. Turning points and sign regimes}

For controlling extrema of \(g\), the most useful coefficient is
\begin{equation*}
H(x)=\kappa-\frac{U(x)^2}{4A}
=\frac{F(x)}{A},
\end{equation*}
where
\begin{equation}
\boxed{
F(x)=\kappa(1-x^2)-\frac14(d-sx)^2.
}
\tag{7}
\end{equation}

Since \(A>0\) on \((-1,1)\), the sign of \(H\) is the sign of \(F\). The polynomial \(F\) is a concave quadratic. Its roots are
\begin{equation}
\boxed{
x_\pm
=
\frac{ds\pm 4\sqrt{\kappa(\kappa+\alpha\beta)}}{s^2+4\kappa}.
}
\tag{8}
\end{equation}
For \(\alpha,\beta>0\), one has
\begin{equation*}
-1<x_-<x_+<1,
\end{equation*}
and
\begin{equation*}
F(x)>0 \quad\Longleftrightarrow\quad x\in(x_-,x_+).
\end{equation*}
Also
\begin{equation*}
F(-1)=-\beta^2,\qquad F(1)=-\alpha^2.
\end{equation*}

The vertex is
\begin{equation}
\boxed{
x_0=\frac{ds}{s^2+4\kappa}.
}
\tag{9}
\end{equation}
Thus \(F\) increases on \((x_-,x_0)\) and decreases on \((x_0,x_+)\).

A useful extremum fact follows immediately from (2). If \(x_*\in(-1,1)\) is a nonzero local extremum of \(g\), then \(g'(x_*)=0\), so
\begin{equation*}
A(x_*)g''(x_*)+H(x_*)g(x_*)=0.
\end{equation*}
At a positive local maximum or negative local minimum, this forces
\begin{equation*}
H(x_*)\ge 0.
\end{equation*}
Hence every nonzero local extremum of (|g|) lies in
\begin{equation*}
\boxed{
[x_-,x_+].
}
\end{equation*}

For the Liouville-normal equation (5), the turning points are the zeros of \(Q_\theta\). In \(x=\cos\theta\), these solve
\begin{equation}
\boxed{
4\kappa(1-x^2)+2-x^2-(d-sx)^2=0.
}
\tag{10}
\end{equation}
The numerator is again concave. Its endpoint values are
\begin{equation*}
1-4\beta^2\quad\text{at }x=-1,
\qquad
1-4\alpha^2\quad\text{at }x=1.
\end{equation*}
Therefore:

\begin{itemize}
\item if \(0\le \alpha,\beta\le \tfrac12\), then \(Q_\theta\ge0\) throughout \((0,\pi)\);
\item if \(\alpha>\tfrac12\), then \(Q_\theta\to-\infty\) as \(\theta\downarrow0\);
\item if \(\beta>\tfrac12\), then \(Q_\theta\to-\infty\) as \(\theta\uparrow\pi\).
\end{itemize}
The \(F\)-turning points are the relevant ones for extrema of \(g\); the \(Q_\theta\)-turning points are the relevant ones for oscillation of the Liouville variable \(Z\).

\section*{3. Candidate Sonin monotonicity functional}

On the oscillatory interval \(F>0\), define
\begin{equation}
\boxed{
\mathcal S_g(x)
=
g(x)^2+\frac{A(x)^2g'(x)^2}{F(x)}.
}
\tag{11}
\end{equation}
Using (2), one obtains the exact Sonin identity
\begin{equation}
\boxed{
\mathcal S_g'(x)
=
-\frac{A(x)^2F'(x)}{F(x)^2},g'(x)^2.
}
\tag{12}
\end{equation}
Here
\begin{equation*}
F'(x)=\frac{ds}{2}-\left(2\kappa+\frac{s^2}{2}\right)x.
\end{equation*}
Hence
\begin{equation*}
F'(x)>0\quad\text{on }(x_-,x_0),
\qquad
F'(x)<0\quad\text{on }(x_0,x_+),
\end{equation*}
so
\begin{equation}
\boxed{
\mathcal S_g \text{ decreases on }(x_-,x_0)
\text{ and increases on }(x_0,x_+).
}
\tag{13}
\end{equation}

At a local extremum of \(g\), \(g'=0\), so
\begin{equation*}
\mathcal S_g(x_*)=g(x_*)^2.
\end{equation*}
Thus the Sonin functional implies that the local extremal values of (|g|) are monotone as one moves from the center \(x_0\) toward either turning point. This reduces the sharp bound to controlling the two outermost lobes adjacent to \(x_-\) and \(x_+\).

For the Liouville equation (5), the parallel functional is
\begin{equation}
\boxed{
\mathcal S_Z(\theta)
=
Z(\theta)^2+\frac{Z'(\theta)^2}{Q_\theta(\theta)}
}
\tag{14}
\end{equation}
on intervals where \(Q_\theta>0\), and
\begin{equation}
\boxed{
\mathcal S_Z'(\theta)
=
-\frac{Q_\theta'(\theta)}{Q_\theta(\theta)^2}Z'(\theta)^2.
}
\tag{15}
\end{equation}
This is useful in the oscillatory bulk, but because \(g=Z/\sqrt{\sin\theta}\), it does not by itself control the endpoint layers.

\section*{4. A closed proof in a nontrivial regime: \(\alpha,\beta\ge 2n,\ n\ge1\)}

Let
\begin{equation*}
R_n(\alpha,\beta)
=
\frac{(n+1)(n+\alpha+\beta+1)}
{(n+\alpha+1)(n+\beta+1)}.
\end{equation*}
The target inequality is
\begin{equation}
|g_n^{(\alpha,\beta)}(x)|\le R_n(\alpha,\beta)^{1/4}.
\tag{16}
\end{equation}

I now prove (16) for
\begin{equation}
\boxed{
n\ge1,\qquad \alpha,\beta\ge 2n.
}
\tag{17}
\end{equation}
This proof uses only the Sturm–Liouville equation and weighted energy identities.

First, for \(\alpha,\beta>0\),
\begin{equation*}
g(-1)=g(1)=0.
\end{equation*}
The following endpoint-weight identities hold:
\begin{equation}
\boxed{
\int_{-1}^{1}\frac{g(x)^2}{1-x},dx=\frac1\alpha,
\qquad
\int_{-1}^{1}\frac{g(x)^2}{1+x},dx=\frac1\beta.
}
\tag{18}
\end{equation}
They follow from the standard connection formula
\begin{equation*}
P_n^{(\alpha,\beta)}
=
\frac{2n+s}{n+s}P_n^{(\alpha-1,\beta)}
+
\frac{n+\beta}{n+s}P_{n-1}^{(\alpha-1,\beta)}
\end{equation*}
and the Jacobi norm formula, with the symmetric identity for lowering \(\beta\).

Consequently,
\begin{equation}
\boxed{
J:=\int_{-1}^{1}\frac{g(x)^2}{1-x^2},dx
=
\frac{s}{2\alpha\beta}.
}
\tag{19}
\end{equation}

Next define the energy
\begin{equation*}
E:=\int_{-1}^{1}A(x)g'(x)^2,dx.
\end{equation*}
Multiplying (2) by \(g\) and integrating by parts gives
\begin{equation*}
E
=
\kappa\int_{-1}^{1}g^2,dx
-\frac14\int_{-1}^{1}\frac{U^2}{A}g^2,dx.
\end{equation*}
Using
\begin{equation*}
\frac{U^2}{A}
=
\frac{2\alpha^2}{1-x}
+
\frac{2\beta^2}{1+x}
-s^2
\end{equation*}
together with (18) and
\begin{equation*}
\int_{-1}^{1}g^2,dx=\frac{2}{2n+s+1},
\end{equation*}
one obtains
\begin{equation}
\boxed{
E
=
n+\frac12-\frac{1}{2(2n+s+1)}.
}
\tag{20}
\end{equation}

Now for any fixed \(x\in[-1,1]\),
\begin{equation*}
g(x)^2
=
2\int_{-1}^{x}g(t)g'(t),dt
\end{equation*}
and also
\begin{equation*}
g(x)^2
=
-2\int_x^1g(t)g'(t),dt.
\end{equation*}
By Cauchy–Schwarz with weights \(A\) and \(A^{-1}\),
\begin{equation*}
g(x)^4
\le
4L(x)P(x),
\qquad
g(x)^4
\le
4M(x)Q(x),
\end{equation*}
where
\begin{equation*}
L(x)=\int_{-1}^{x}\frac{g^2}{A},\quad
M(x)=\int_x^1\frac{g^2}{A},
\end{equation*}
\begin{equation*}
P(x)=\int_{-1}^{x}Ag'^2,\quad
Q(x)=\int_x^1Ag'^2.
\end{equation*}
Since \(L+M=J\) and \(P+Q=E\),
\begin{equation*}
\min{L P,M Q}\le \frac{JE}{4}.
\end{equation*}
Therefore
\begin{equation}
\boxed{
|g(x)|^4\le JE
=
\frac{s}{2\alpha\beta}
\left(
n+\frac12-\frac{1}{2(2n+s+1)}
\right).
}
\tag{21}
\end{equation}

It remains to show that \(JE\le R_n(\alpha,\beta)\) under \(\alpha,\beta\ge2n\).

The inequality \(JE\le R_n\) is equivalent to
\begin{equation}
\frac{E}{2},
\frac{s}{n+s+1}
\left(1+\frac{n+1}{\alpha}\right)
\left(1+\frac{n+1}{\beta}\right)
\le n+1.
\tag{22}
\end{equation}
Let
\begin{equation*}
T(\alpha,\beta)
=
\frac{s}{n+s+1}
\left(1+\frac{n+1}{\alpha}\right)
\left(1+\frac{n+1}{\beta}\right).
\end{equation*}
For fixed \(s\), the product is maximized when \(\alpha\beta\) is minimized. Under \(\alpha,\beta\ge2n\), this occurs at one endpoint, say \(\alpha=2n\), \(\beta=s-2n\). Hence
\begin{equation*}
T(\alpha,\beta)
\le
\frac{s}{n+s+1}
\left(1+\frac{n+1}{2n}\right)
\left(1+\frac{n+1}{s-2n}\right).
\end{equation*}
The right-hand side is decreasing for \(s\ge4n\), because its derivative equals
\begin{equation*}
-\frac{(n+1)(3n+1)(2s-n+1)}
{(s-2n)^2(n+s+1)^2}<0.
\end{equation*}
Thus its maximum occurs at \(s=4n\), giving
\begin{equation*}
T(\alpha,\beta)
\le
\frac{4n}{5n+1}
\left(1+\frac{n+1}{2n}\right)^2
=
\frac{(3n+1)^2}{n(5n+1)}.
\end{equation*}
Finally,
\begin{equation*}
\frac{(3n+1)^2}{n(5n+1)}
\le
\frac{4(n+1)}{2n+1},
\end{equation*}
because the difference after cross-multiplication is
\begin{equation*}
(n-1)(2n^2+5n+1)\ge0.
\end{equation*}
Since
\begin{equation*}
E<n+\frac12=\frac{2n+1}{2},
\end{equation*}
we get
\begin{equation*}
\frac{E}{2}T(\alpha,\beta)
<
\frac{2n+1}{4}\cdot
\frac{4(n+1)}{2n+1}
=
n+1.
\end{equation*}
So \(JE\le R_n\), and (21) yields
\begin{equation*}
|g(x)|^4\le R_n(\alpha,\beta).
\end{equation*}
Therefore
\begin{equation*}
\boxed{
|g_n^{(\alpha,\beta)}(x)|
\le
\left(
\frac{(n+1)(n+\alpha+\beta+1)}
{(n+\alpha+1)(n+\beta+1)}
\right)^{1/4}
}
\end{equation*}
for all \(x\in[-1,1]\), \(n\ge1\), and \(\alpha,\beta\ge2n\).

No integer-parameter representation theory is used.

\section*{5. What remains}

The Sonin functional (11) gives a good structural reduction, but it does not by itself close the sharp KKT bound. The unresolved analytic work is concentrated in these regimes:

\begin{enumerate}
\item \textbf{Endpoint-turning lobes.}
Since \(\mathcal S_g\) is monotone away from the center \(x_0\), the largest local extrema are forced toward the outer lobes adjacent to \(x_-\) and \(x_+\). Sharp control there requires Bessel/Airy-type endpoint estimates.
\item \textbf{Small parameter edge:}
\begin{equation*}
0<\alpha<2n
\quad\text{or}\quad
0<\beta<2n.
\end{equation*}
This includes the most delicate limits \(\alpha\downarrow0\) or \(\beta\downarrow0\), where the target constant tends to (1) and the maximum migrates toward an endpoint.
\item \textbf{Strongly unbalanced parameters:}
One of \(\alpha,\beta\) small and the other large. Then one direct turning point approaches an endpoint, and the simple energy bound above is too weak.
\item \textbf{Liouville endpoint thresholds:}
In the normal form (5), the sign of the endpoint inverse-square term changes at
\begin{equation*}
\alpha=\frac12,\qquad \beta=\frac12.
\end{equation*}
These thresholds likely require separate Bessel-regime estimates.
\item \textbf{Finite small-\(n\) algebraic cases.}
Cases such as \(n=1,2\) can be attacked by explicit polynomial maximization in the beta variable \(t=(1+x)/2\). They are good test cases for identifying the missing endpoint inequalities.
\end{enumerate}
The exact ODE and Sonin identity are strong enough to reduce the global problem to endpoint/turning-point majorants. The proof above shows that the same Sturm–Liouville framework already proves the target inequality in the large-balanced regime \(\alpha,\beta\ge2n\).
\end{document}
